% NAF 4073, batch 1 — Gallica views 1–20.
% Pass: Opus 5 (claude-opus-5), 2026-09-11 — first pass, unchecked against
% the views by a human.
%
% What the twenty views hold. NAF 4073 is the BnF's binding of letters
% addressed to Sophie Germain, from the papers of Libri; the leaves are
% mounted singly in-4°, and most views photograph two mounted leaves at once.
% This batch is the binding and the first thirteen folios: not one line is in
% Germain's hand — every letter here is one she received. All are French
% prose; there is no mathematics set as formulae, but four of the six letters
% are about her elasticity work and one is about Gauss.
%
%   v. 6      f. 1r     A. Choron to Germain, « Paris le 29. Janvier 1818 »,
%                       sending her a copy of his « observations sur
%                       l'enseignement simultané de la lecture et de
%                       l'écriture en musique ». Choron is NOT among the
%                       correspondents the BnF's notice lists.
%   v. 7–8    f. 2r–2v  Delambre, on the letterhead « Université Impériale /
%                       Trésorerie », « Paris, ce lundi 14 mai 1810 »: Gauss,
%                       awarded the Lalande medal, would rather have a clock
%                       than the 380 f. remaining in silver, and Delambre asks
%                       Germain to choose it. He quotes Gauss's own words and
%                       corrects his German (« Pendeluhr » is a pendulum
%                       clock, not a watch).
%   v. 9      f. 3v     The address leaf of that letter: « maison de
%                       M. Doulcet d'Egligny maire du 7e arrondissement, Rue
%                       Ste Croix de la bretonnerie ».
%   v. 10–11  f. 4r–4v  Delambre, « Paris le 21 juillet 1821 », answering a
%                       letter and the « savant mémoire » that came with it
%                       — she is combating the principles of a « Géomètre
%                       célèbre » — then explaining at length why she cannot
%                       be seated at the Académie's public séances: the
%                       « billets du centre », forty tickets for
%                       seventy-five academicians.
%   v. 12     f. 5v     The address leaf of that letter: « Rue de Braque au
%                       marais, N°. 47 ou 41 » — Delambre unsure of the number.
%   v. 12     f. 6r     Fourier, « Paris 29 Sebre 1820 », pressing her to
%                       publish the theory she is working on: « personne ne
%                       peut traiter avec plus de succès cette difficile et
%                       intéressante question ».
%   v. 14     f. 7v     The address leaf of that letter, giving the number:
%                       « Rue de Braque au marais N° 41 ».
%   v. 14     f. 8r     Fourier, « Paris 30 janvier 1827 »: he is too ill to
%                       come, will send his éloge de Delambre, and defends the
%                       Académie française's démarche « objet d'une colère si
%                       injuste et si violente », the Académie des sciences
%                       having not joined it.
%   v. 16     f. 9v     An address leaf without a street: hand-delivered.
%   v. 16     f. 10r    Fourier, « vendredi matin », undated: someone in his
%                       house is ill; Libri has written; he will report on
%                       Cauchy's visit.
%   v. 18     f. 11v    A second address leaf without a street.
%   v. 19     f. 12r    Fourier, « mardi soir », undated: he sends the number
%                       of the Annales de physique carrying Savart's memoir on
%                       the vibrations of elastic plates, and promises Navier's
%                       memoir on the flexion of elastic plates, which Navier
%                       had lithographed. Greetings to « madame sa mère », who
%                       died in 1823.
%
% Absent from the output, holding nothing of hers or of theirs: v. 1 (front
% cover), v. 2 (pastedown and the shelfmark label), v. 3 (the guard leaf's
% library slips, « Volume de 32 feuillets, 15 Mars 1875 »), v. 4 (a later
% cataloguer's pencilled contents list, « Lettres de et à Sophie Germain —
% Fourier … Delambre — Lalande — De la Grange »), v. 5, 13, 15, 17, 20 (blank
% versos on which the facing letter shows through, mirrored). The BnF's stamps
% and the wax seals are not transcribed.
%
% Folios are the pencilled numbers on the rectos (1–13 across the batch); a
% verso is written as the recto's number with v, the leaf being the same. On
% several views the address shows through the paper of the facing leaf,
% mirrored; that show-through is never transcribed — the address is given only
% at the view where it is right-reading.
%
% Notation set once and held to. The abbreviations « Mr », « M. », « Ste »,
% « 3me », « 7e » stay as written. « n° » is set with the degree sign. In the
% 1810 letter the sum is written with a superscript franc sign after the
% figure; it is set here as « 380 f. ». Fourier writes septembre as an « S »
% with a superscript « bre »; it is set on the line as « Sebre ». Words
% Delambre or Fourier underlined are set with \emph. A heavy stroke that
% cancels a word is \struck{}; where the cancelled word cannot be read, \ill{}
% stands inside a \note{}.
\documentclass[11pt,a4paper]{article}
% The preamble shared by every transcription of the mathematicians' manuscripts in Gallica.
%
% It rests on one idea: **uncertainty compounds, it does not resolve.** A
% transcription that smooths over a doubtful word has destroyed the only thing
% it was made for — a reader must be able to tell, at every point, what was on
% the page and what was supplied. The apparatus macros below are therefore the
% critical apparatus, not decoration.
%
% The same commands are recognised by `scripts/render.mjs`, which produces the
% site's reading view, and by `scripts/tei.mjs`. A macro added here must be
% added there too, or rendering fails loudly — which is the wanted behaviour.
%
% Comments here are English, like the rest of the repository. The documents
% this preamble sets are French, and that is deliberate: see the skills.

\ProvidesPackage{germain}[2026/09/15 Transcriptions of mathematicians' manuscripts in Gallica]

\usepackage{amsmath,amssymb,amsthm}
\usepackage{mathrsfs}
% Kept from the parent preamble so the renderer's diagram support stays
% available; these manuscripts draw no commutative diagrams, and nothing here uses it.
\usepackage{tikz-cd}
\usepackage[english,french]{babel}
\usepackage{fontspec}

% --- Greek ------------------------------------------------------------------
% The text face stays fontspec's default, Latin Modern, which has no Greek:
% XeTeX drops every glyph it lacks with a line in the log and nothing on the
% page, and the Greek words the manuscripts quote vanished from the PDFs. So
% the Greek and Greek Extended blocks get a XeTeX character class of their own,
% and entering or leaving a run of it switches the family to CMU Serif — the
% same Computer Modern design, with polytonic Greek — and back. Only the
% family changes: a Greek word in \textit or \textbf keeps its shape and series.
% The transitions to and from every other class carry the tokens that class
% has towards an ordinary letter, so French spacing before « ; » or inside
% « » still holds after a Greek word. No grouping, so no brace or \endgroup
% can come out unbalanced; maths is left alone, since XeTeX inserts nothing
% there. Kept identical in `exercices.sty`.
\makeatletter
\newfontfamily\germainGreekfont{cmun}[
  Extension=.otf, NFSSFamily=germaingreek,
  UprightFont=*rm, ItalicFont=*ti, SlantedFont=*sl,
  BoldFont=*bx, BoldItalicFont=*bi, BoldSlantedFont=*bl]
\newXeTeXintercharclass\germain@greek
\def\germain@greekrange#1#2{%
  \@tempcnta=#1\relax
  \loop
    \XeTeXcharclass\@tempcnta=\germain@greek
    \ifnum\@tempcnta<#2\advance\@tempcnta\@ne\repeat}
\germain@greekrange{"0370}{"03FF}
\germain@greekrange{"1F00}{"1FFF}
\def\germain@greekprev{\rmdefault}
\def\germain@greekin{%
  \edef\germain@greekprev{\f@family}\fontfamily{germaingreek}\selectfont}
\def\germain@greekout{\fontfamily{\germain@greekprev}\selectfont}
\def\germain@greekjoin{%
  \XeTeXinterchartokenstate=\@ne
  \@tempcnta=\z@
  \loop
    \ifnum\@tempcnta=\germain@greek\else
      \XeTeXinterchartoks\germain@greek\@tempcnta=\expandafter{%
        \expandafter\germain@greekout\the\XeTeXinterchartoks\z@\@tempcnta}%
      \XeTeXinterchartoks\@tempcnta\germain@greek=\expandafter{%
        \the\XeTeXinterchartoks\@tempcnta\z@\germain@greekin}%
    \fi
    \ifnum\@tempcnta<4095\advance\@tempcnta\@ne\repeat}
% babel-french sets its own classes, and saves and restores the interchar
% state, each time a language is selected: join again after it does.
\addto\extrasfrench{\germain@greekjoin}
\addto\extrasenglish{\germain@greekjoin}
\AtBeginDocument{\germain@greekjoin}
\makeatother

\usepackage{geometry}
\geometry{a4paper,margin=25mm,marginparwidth=28mm}
\usepackage{marginnote}
\usepackage{xcolor}
\usepackage[normalem]{ulem}
\setlength{\emergencystretch}{3em}
\usepackage{hyperref}
\hypersetup{colorlinks=true,linkcolor=black,urlcolor=black,pdfborder={0 0 0}}

% --- Batch metadata ---------------------------------------------------------
% Taken as they stand by the rendering script, which shows them at the head of
% the reading view. They fix what the file claims to cover. The macro names are
% the parent project's, so that the scripts stay shared; \folder here means the
% volume's slug — `fr-9115` — and \pages counts Gallica views.

\newcommand{\folder}[1]{\gdef\grothFolder{#1}}
\newcommand{\batch}[1]{\gdef\grothBatch{#1}}
\newcommand{\pages}[2]{\gdef\grothFirst{#1}\gdef\grothLast{#2}}
\newcommand{\foldertitle}[1]{\gdef\grothFoldertitle{#1}}

% The holder's own shelfmark, printed as they write it: « Français 9115 ».
\newcommand{\shelfmark}[1]{\gdef\grothShelfmark{#1}}

% The dating, copied verbatim from the holder's catalogue — never inferred
% here. « XIXe siècle » is the BnF's; a pass that dates a leaf from its content
% says so in a \note{}, as its own inference.
\newcommand{\dating}[1]{\gdef\grothDating{#1}}

% The status notice, printed in the title block and repeated as a diagonal
% watermark on every page of the PDF. The manuscripts are in the public
% domain; what this line declares is not a right but a quality — an
% unverified first pass by a machine. The skills write the line; a file
% without it gets no watermark, so leaving it out is visible in the output.
\newcommand{\watermark}[1]{\gdef\grothWatermark{#1}}

\gdef\grothFolder{?}\gdef\grothBatch{}\gdef\grothFirst{?}\gdef\grothLast{?}%
\gdef\grothFoldertitle{}\gdef\grothShelfmark{}\gdef\grothDating{}\gdef\grothWatermark{}

% --- The résumé -------------------------------------------------------------
\newenvironment{resume}
  {\small\setlength{\parskip}{0.45em}}
  {\par\medskip\hrule\medskip}

% --- Keywords ---------------------------------------------------------------
% In English, closing the modernised reading's résumé: the modern vocabulary
% under which to search for what the leaves construct. The manifest extracts
% them to tag the volume — this line is the single source of the tags.
\newcommand{\keywords}[1]{\par\smallskip\noindent
  {\footnotesize\textsc{Keywords} --- \textit{#1}}\par}

% --- The critical apparatus -------------------------------------------------

% The Gallica view number, in the margin. This is the anchor that lets a
% disagreement over a formula be settled by looking at *one* image rather than
% a volume of 750. The site uses it to turn the facsimile as the transcription
% is scrolled. A view is one scanned image: usually one side of a leaf.
\newcommand{\page}[1]{%
  \marginnote{\footnotesize\color{gray!70}\textsf{v.\,#1}}%
  \label{page:#1}\ignorespaces}

% The BnF's pencilled foliation, where the leaf carries it and a pass can read
% it: « 198r ». This is what the literature cites — « ff. 198r–208v » — and
% the one line that turns a citation into a view. Recorded, never inferred:
% a folio a pass cannot read is not written.
\newcommand{\folio}[1]{%
  \marginnote{\footnotesize\color{gray!50}\textsf{f.\,#1}}\ignorespaces}

% The modernised reading's coarser counterpart to \page{}: which views a
% section is drawn from.
\newcommand{\pagerange}[2]{%
  \marginnote{\footnotesize\color{gray!70}\textsf{v.\,#1--#2}}%
  \ignorespaces}

% Illegible: never guessed.
\newcommand{\ill}{\textcolor{red!60!black}{[\,\ldots\,]}}

% A doubtful reading: the word is given, and flagged as uncertain.
\newcommand{\uncertain}[1]{\uline{#1}}

% An editorial addition — a missing letter, an implied word.
\newcommand{\add}[1]{\textcolor{blue!50!black}{[#1]}}

% Struck out by the writer. What was crossed out is part of the document.
\newcommand{\struck}[1]{\sout{#1}}

% The transcriber's note, on the state of the page or the mathematics.
\newcommand{\note}[1]{\par\smallskip{\small\color{gray!60!black}\textit{[#1]}}\par\smallskip}

% A marginal note of hers, distinct from the body of the text.
\newcommand{\marginal}[1]{\par\smallskip\noindent\hspace{1em}%
  {\small\color{gray!60!black}#1}\par\smallskip}

% --- Title ------------------------------------------------------------------
% The title block is French, like the documents themselves.

\AddToHook{shipout/background}{%
  \ifx\grothWatermark\empty\else
    \begin{tikzpicture}[remember picture, overlay]
      \node[rotate=52, scale=3.4, align=center, text=gray!18,
            font=\sffamily\bfseries]
        at (current page.center) {\grothWatermark};
    \end{tikzpicture}%
  \fi}

\AtBeginDocument{%
  \begin{center}
    {\large\bfseries Manuscrits de mathématiciens (Gallica) ---
      \ifx\grothShelfmark\empty\grothFolder\else\grothShelfmark\fi}\\[2pt]
    {\itshape\grothFoldertitle}\\[4pt]
    {\small vues \grothFirst--\grothLast
      \ifx\grothBatch\empty\else\ (lot \grothBatch)\fi}\\[2pt]
    \ifx\grothDating\empty\else
      {\small\color{gray!60!black}Datation du catalogue : \grothDating}\\[2pt]
    \fi
    \ifx\grothWatermark\empty\else
      {\small\bfseries\color{red!45!gray}\grothWatermark}\\[2pt]
    \fi
    {\footnotesize\color{gray!60!black}Transcription établie d'après les images IIIF de Gallica.
     Source gallica.bnf.fr / Bibliothèque nationale de France.\\
     Les lectures incertaines sont soulignées, les illisibles notées [\,\ldots\,],
     les ajouts entre crochets ; \textsf{v.} numérote les vues de Gallica,
     \textsf{f.} la foliotation au crayon de la BnF.}
  \end{center}
  \vspace{1em}}

\endinput


\folder{naf-4073}
\shelfmark{NAF 4073}
\batch{1}
\pages{1}{20}
\dating{XVIIIe siècle}
\watermark{Lecture automatique\\non vérifiée}
\foldertitle{Correspondance de Sophie GERMAIN.}

\begin{document}

\section{Choron à Sophie Germain, 29 janvier 1818}

\page{6}\folio{1r}
J'ai l'honneur d'offrir mes hommages à Mlle Germain et de lui adresser
ci-inclus l'exemplaire qu'elle m'a permis de lui présenter de mes
observations sur l'enseignement simultané de la lecture et de l'écriture en
musique. Je désire qu'elle en soit satisfaite

A. Choron

Paris le 29. Janvier 1818.

\note{le nom de Choron ne figure pas dans la notice de la Bibliothèque, qui
énumère Delambre, Fourier, Germain, Lagrange, Lalande et Libri}

\section{Delambre à Sophie Germain, 14 mai 1810 : la médaille de Lalande et la pendule de Gauss}

\page{7}\folio{2r}
Université Impériale.

Trésorerie.

\note{les deux lignes sont en tête de la feuille, d'une écriture posée qui
peut être celle d'un en-tête gravé}

Paris, ce lundi 14 mai 1810.

Par une lettre que je viens de recevoir de Mr Gauss je suis chargé d'une
commission sur laquelle il m'engage à prendre vos avis. dans toute autre
circonstance je n'eusse pas manqué de saisir avec empressement cette occasion
pour vous porter mes hommages. Je viens de rendre les derniers devoirs à la
mère de ma femme. ce triste événement, joint à toutes mes autres occupations,
ne me laissera de quelque tems aucun moment dont je puisse disposer.

sur la somme de 500 fr. valeur de la médaille fondée par Lalande et qui vient
de lui être décernée Mr Gauss desire avoir une pendule ; voici les termes de
sa lettre

\begin{quote}
au lieu d'accepter le reste savoir 380 f. en argent j'aimerois mieux avoir
\emph{une belle montre à Pendule}. je n'en fixe pas le prix ; qu'il soit de
60 ou de 300 fr. cela m'est indifférent pourvû que la montre soit assez
élégante pour être offerte en cadeau à mon épouse et pour pouvoir servir de
décoration à sa chambre
\end{quote}

\page{8}\folio{2v}
\begin{quote}
peut-être Mademoiselle Sophie Germain ( à laquelle je vous prie de faire
mille complimens de ma part ) auroit la bonté de se charger du choix.
\end{quote}

Il me paroit que c'est une Pendule et non une montre que désire Mr Gauss. il
a traduit par \emph{montre à Pendule} l'expression allemande Pendeluhr. il
s'agit donc de lui choisir une Pendule du prix de 60 à 300 fr. mais il veut
qu'elle soit \emph{élégante} et je craindrois de ne pas deviner bien juste le
goût de Madame Gauss. Veuillez donc, Mademoiselle, m'informer si je puis me
flatter d'être aidé de vos conseils, et me faire parvenir vos ordres.
J'espère avoir plus de liberté dans quelques jours et j'en profiterois pour
aller apprendre le résultat de vos idées ou de vos soins et je m'occuperois
aussitôt après de faire partir la pendule pour Gottingue en profitant, s'il
en est tems encore, des occasions qu'il m'a indiquées.

Agréez l'hommage des sentimens respectueux avec lesquels j'ai l'honneur
d'être.

Mademoiselle

Votre très humble et très obéissant serviteur

Delambre

Palais du Corps législatif

\note{la feuille voisine, f. 3r, ne porte rien : on y voit seulement, à
l'envers, l'adresse écrite sur son autre face}

\page{9}\folio{3v}
A Mademoiselle

Mademoiselle Sophie Germain

maison de M. Doulcet d'Egligny maire

du 7e arrondissement

Rue Ste Croix de la bretonnerie. n°. \uncertain{25}

Paris.

\note{le dernier chiffre est chargé d'encre ; on lit 25 plutôt que 23}

\section{Delambre à Sophie Germain, 21 juillet 1821 : les billets du centre}

\page{10}\folio{4r}
Mademoiselle

Avant de répondre à la lettre dont vous m'avez honoré, j'ai cédé à
l'empressement de lire le savant mémoire qui l'accompagnoit et dont je
n'avois vû que le titre, lorsque je l'avois annoncé à l'académie à sa
dernière séance. il ne m'appartient pas et peut-être n'appartient-il à
personne de prononcer sur les questions que vous regardez comme douteuses.
mais ce qui obtiendra une approbation générale c'est la réserve avec laquelle
vous exposez vos opinions et vos conjectures, ce sont les égards avec
lesquels vous traitez le Géomètre célèbre dont vous combattez les principes
qui vous ont paru ne pouvoir s'accorder avec ce que vous démontrez.

\note{le mémoire est vraisemblablement les « Recherches sur la théorie des
surfaces élastiques », parues en 1821 ; le rapprochement est celui de la
présente lecture et non du document}

L'ordre établi depuis plusieurs années dans la distribution des places dans
les séances publiques me semble comme à vous susceptible de plus d'une
objection, l'auteur en est feu Mr Suard, et il n'est pas un membre de
l'Institut, surtout s'il est marié, qui n'ait eu plus d'une fois, à s'en
plaindre. Mais depuis l'ordonnance royale chaque académie peut disposer du
local suivant ses convenances. l'académie française s'étant emparée de toutes
les places de choix, aux jours où elle préside, chacune des autres académies
a voulu probablement et le privilégié à son tour. on a conservé l'usage des
\emph{billets du centre}. ils sont répartis parmi nous entre le Président, le
Vice-président, les deux secrétaires perpétuels et les lecteurs. il est un
nombre de ces billets qu'on réserve pour les grands fonctionnaires et pour
les étrangers célèbres, en sorte qu'à chaque séance de notre académie j'ai
tout au plus dix de ces billets. je me suis fait une loi de les distribuer à
ceux de mes confrères, qui m'en demandent pour leurs femmes

\page{11}\folio{4v}
Le Bureau n'a pas plus de 40 billets à distribuer et les académiciens sont au
nombre de 75. heureusement leurs femmes ne montrent pas un grand
empressement en sorte que jusqu'ici je n'ai pas eu le chagrin d'en refuser
une seule. c'est aux Perpétuels et non à ceux qui changent chaque année que
les étrangers s'adressent de préférence.

\note{un pâté d'encre couvre un mot après « une seule » ; il n'a pu être lu}

mais si j'avois été instruit du désir que vous aviez d'assister à nos séances
publiques, j'aurois trouvé le moyen de vous réserver un de ces billets et je
m'en souviendrai certainement au mois de Mars prochain. la chose ne seroit
pas impossible au mois de juillet. le secrétaire perpétuel de l'académie des
inscriptions me gratifie ordinairement de plusieurs billets du centre. mais
pour les deux autres académies, je n'en ai pas un seul et si ma femme n'avait
pas renoncé depuis longtems à ces solemnités je serois obligé de m'adresser
pour elle au secrétaire perpétuel en fonction. je ne m'y suis hasardé qu'une
seule fois et ma demande a été accueillie ; j'aime à croire qu'elle le seroit
toutes les fois que je m'y prendrois à tems. j'en courrai volontiers le
hasard toutes les fois que vous désirerez un de ces billets et que vous m'en
aurez instruit au plus tard le lundi qui précédera la séance.

Agréez l'hommage des sentimens respectueux avec lesquels j'ai l'honneur
d'être.

Mademoiselle

Votre très humble et très obéissant serviteur

Delambre.

Paris le 21 juillet 1821

\page{12}\folio{5v}
A Mademoiselle

Mademoiselle Sophie Germain

Rue de Braque au marais

N°. 47 ou 41

\note{l'adresse est écrite en travers de la feuille ; Delambre hésite sur le
numéro, que l'adresse du f. 7v donnera comme 41}

\section{Fourier à Sophie Germain, 29 septembre 1820 : « je vous invite instamment à publier »}

\page{12}\folio{6r}
Paris 29 Sebre 1820

Mademoiselle,

Je regrette extrêmement de partir pour la campagne, sans avoir eu l'honneur
de vous voir ; j'y demeurerai quinze jours environ, et à mon retour je me
présenterai chez vous, si vous voulez bien me le permettre. Je suis persuadé
que vous aurez donné de nouveaux développemens à la théorie dont vous vous
occupez, et que je vous invite instamment à publier. personne ne peut traiter
avec plus de succès cette difficile et intéressante question.

agréez mademoiselle l'hommage de mon respect.

J. Fourier

\section{Fourier à Sophie Germain, 30 janvier 1827 : l'éloge de Delambre et la colère contre l'Académie française}

\page{14}\folio{7v}
A Mademoiselle

Mademoiselle Sophie Germain

Rue de Braque au marais N° 41

Paris.

\page{14}\folio{8r}
Mademoiselle

combien je regrette de ne pouvoir profiter de l'invitation que vous avez la
bonté de m'offrir. vous connaîtrez les raisons de santé qui m'obligent de
renoncer à tout ce qui pourroit m'être le plus agréable.

\note{un mot est biffé après « je regrette » ; il n'a pu être lu}

J'aurai l'honneur de vous envoyer ou de vous porter moi même un exemplaire de
l'éloge de Delambre. je vois avec le plus grand plaisir que vous porterez
votre attention sur ces considérations nouvelles que présente le calcul des
inégalités.

\note{un trait appuyé traverse le mot « éloge » sans qu'on puisse dire s'il
biffe ou souligne}

j'ai \ill{} je \uncertain{rends} tems à l'académie un \ill{} sur divers points
d'analyse qui intéressent la théorie de la chaleur et le mouvement
\uncertain{diurne} de la chaleur terrestre.

\note{deux mots de cette phrase sont biffés et n'ont pu être lus ; la syntaxe
de la phrase n'a pu être rétablie}

on auroit beaucoup \ill{} été \ill{} l'académie des sciences à suivre une
démarche semblable à celle qui vient d'être l'objet d'une colère si injuste
et si violente. mais nous n'avons pas cru devoir nous occuper de \ill{}
matière si \uncertain{salutaire}, quoique très frappés de l'entreprise qui
menace les lettres.

\note{un pâté couvre un mot après « nous occuper de » ; deux autres mots de
la première phrase sont illisibles. La démarche dont il s'agit est
vraisemblablement l'adresse de l'Académie française contre le projet de loi
sur la presse, en janvier 1827 ; le rapprochement est celui de la présente
lecture}

agréez mademoiselle l'hommage de mon respect.

Jh. Fourier

Paris 30 janvier 1827

\section{Fourier à Sophie Germain, « vendredi matin » : Libri et Cauchy}

\page{16}\folio{9v}
Mademoiselle

Mademoiselle Sophie Germain.

\note{ni rue ni ville : le billet a été porté}

\page{16}\folio{10r}
mademoiselle,

Je suis bien reconnoissant de votre souvenir et de l'intérêt que vous inspire
une personne qui souffre auprès de moi. je pense qu'il y a maintenant espoir
de guérison. pardonnez-moi de ne pas me rendre à l'invitation que vous avez
la bonté de me proposer, j'espère être plus heureux une autre fois.

je m'acquitterai de la recommandation que vous me \uncertain{confiez} —
\uncertain{me mande} Mr Libri. il m'en avoit aussi parlé dans sa dernière
lettre. je ferai mon possible pour que ce vœu ne soit pas différé. il nous a
envoyé un nouveau mémoire sur la théorie de la chaleur. je vois avec beaucoup
de plaisir qu'un \uncertain{homme} de son mérite s'occupe de ces questions.

\note{la fin de la première phrase est d'une lecture douteuse : on lit
« que vous me confiez » puis, après un tiret, « me mande Mr Libri », sans que
la construction se laisse rétablir}

j'aurai l'honneur de vous voir dans la semaine prochaine et je vous rendrai
compte de la visite de Mr Cauchy. Mercredi ou jeudi, un autre jour si vous le
préférez.

agréez mademoiselle l'hommage de ma vive reconnaissance et de mon respect.

Fourier.

vendredi matin.

\page{18}\folio{11v}
Mademoiselle

Mademoiselle Sophie Germain.

\section{Fourier à Sophie Germain, « mardi soir » : le mémoire de Savart et celui de Navier}

\page{19}\folio{12r}
j'ai l'honneur d'adresser à mademoiselle germain le dernier numéro des
annales de physique où se trouve le mémoire de Mr Savart sur les vibrations
des lames élastiques. j'ai pensé que la lecture de cet écrit pouvoit
l'intéresser. j'espère lui envoyer demain un mémoire qui a été présenté à
l'académie par Mr Navier et qui a pour objet l'analyse des flexions des plans
élastiques. l'auteur a fait copier ce manuscrit par les presses
lithographiques et m'en a remis un exemplaire.

je prie mademoiselle germain d'agréer l'hommage de mon respect et de vouloir
bien l'offrir de ma part à madame sa mère.

Fourier.

mardi soir.

\note{la mère de Sophie Germain étant morte en 1823, ce billet lui est
antérieur ; la remarque est de la présente lecture et non du document}

\end{document}
